\textcolor{red}{[Topic]} Vision-language models (VLMs) such as CLIP, SigLIP, and ALIGN, along with self-supervised visual encoders like DINOv2, are routinely adapted to new domains using lightweight mechanisms such as LoRA, DoRA, MaPLe, and TAP, that deliver strong empirical gains across classification, segmentation, and retrieval. \textcolor{red}{[Motivation]} While these adapters are widely deployed, the internal changes they induce remain poorly understood, and existing mechanistic-interpretability analyses overwhelmingly rely on a single \emph{generic} sparse autoencoder (SAE) trained once on the unadapted model and reused to explain every adapted variant, an assumption that has never been stress-tested. \textcolor{red}{[Contribution]} In this work, we show across four encoders, four adapter families, and multiple SAE variants that generic SAEs do not transfer faithfully to adapted models: their interpretability degrades substantially under adaptation, with the gap growing non-linearly in the distance between pretraining and target domains, a phenomenon we quantify using a new \emph{Domain-Aware Monosemanticity Score} that measures concept coherence relative to the target distribution rather than the pretraining one. \textcolor{red}{[Detail/Nuance]} Specifically, on datasets ranging from ImageNet-like benchmarks to far-OOD medical (MedMNIST, ChestMNIST) and satellite (EuroSAT) imagery, generic-SAE features lose both reconstruction fidelity and concept selectivity (measured via sparsity, label entropy, and our Domain-Aware Monosemanticity Score) when applied to adapted CLIP, SigLIP, ALIGN, and DINOv2. \textcolor{red}{[Evidence / Contribution 2]} Leveraging this observation, we show that \emph{specific} SAEs, explicitly retrained on the adapted model, recover and surpass the original interpretability profile, yielding sparser and more monosemantic features whose gain over generic SAEs scales with domain distance. \textcolor{red}{[Weaker result]} We further find that adapter families differ mechanistically: prompt-based methods primarily steer existing features, while weight-space adaptation more aggressively reshapes the feature dictionary itself. \textcolor{red}{[Narrow impact]} Our findings underscore the brittleness of current SAE-based interpretation pipelines that reuse generic dictionaries on adapted models. \textcolor{red}{[Broad impact]} More broadly, our work argues that trustworthy mechanistic interpretation of adapted foundation models requires specific SAEs constructed on the adapted model itself, and provides a multi-axis                                            , anchored by the Domain-Aware Monosemanticity Score, to quantify when generic-SAE reuse is justified.